\section{Goals}

\subsection{Practical Significance}
Learning Management Systems (LMS) such as Moodle are widely adopted by universities to organize courses, assign coursework, administer assessments, and manage grades. However, these systems primarily serve as platforms for storing and managing data and do not yet provide intelligent support to learners and instructors. This project aims to add such capabilities on top of the existing Moodle platform, rather than building a new LMS from scratch.
For students, the system allows questions about course content to be posed in natural language and returns answers grounded in the course's own materials, accompanied by source citations. This approach reduces the time students spend searching through large collections of slides, textbooks, and syllabi, and eases the burden on instructors of repeatedly answering the same questions. For instructors, the system provides a tool for the early identification of students at risk of poor academic performance, together with the specific reasons behind each assessment, so that timely reminders and support can be offered instead of problems being discovered only after final grades are released. At the institutional level, the resulting analyses serve as a useful data source for monitoring learning progress and improving the quality of teaching.

\subsection{Scientific and Technical Significance}
From a technical perspective, the project applies Retrieval-Augmented Generation (RAG) to question answering over academic documents in Vietnamese. This task involves several related challenges, including processing documents in multiple formats, combining semantic and keyword-based retrieval, and controlling the tendency of language models to produce answers that deviate from the source material. In parallel, the project applies machine learning to predict academic risk from students' activity data on the LMS, with particular emphasis on the interpretability of predictions so that non-technical users can understand and trust the results.
In addition, building a system composed of multiple components, namely the LMS, backend services, data-processing services, artificial intelligence services, and a user interface, provides an opportunity to practice the entire software development process, from analysis and design through deployment and testing.

\subsection{Significance for the Project Team}
For the project team, this work offers an opportunity to develop teamwork skills within a software engineering process, to practice integrating an existing open-source system with extension components, and to become familiar with artificial intelligence technologies from an applied perspective, aimed at real products rather than purely theoretical research.

\subsection{General Objective}
The general objective of this project is to develop an intelligent learning management system for the university environment, built on the Moodle platform, which integrates artificial intelligence features to support students in their learning and to help instructors monitor and identify early those students at risk of poor academic performance.

\subsection{Specific Objectives}
In terms of business analysis, the team surveys the needs of students, instructors, and administrators in the learning and training management process, and analyzes similar systems such as Moodle, Canvas, and Google Classroom to identify their strengths, weaknesses, and missing functionalities, which form the basis for proposing additional features. On this basis, the functional and non-functional requirements of the system are determined.
In terms of technology and design, the team selects appropriate technologies for each component and designs the overall architecture, database, and user interface in a manner that avoids modifying the Moodle core, in order to ensure maintainability and ease of future upgrades.
In terms of implementation, the team develops a mechanism for synchronizing data from Moodle, including class lists, students, learning materials, grades, and activity logs, to the extension services. Building on this data, the team develops a chatbot that answers questions based on course materials, with enrollment-based access control and source citations, together with an academic early-warning module that produces a list of at-risk students along with the reasons for each. These two features are delivered to users through interfaces designed for students and instructors.
In terms of testing and evaluation, the team verifies the functionality and integration of the system, evaluates chatbot quality in terms of retrieval accuracy and latency, and assesses the quality of the warning model using precision, recall, and F1-score, in order to obtain quantitative evidence on which to base conclusions.